\section{Power tips}
\label{sec:tips}

These are the changes that take a working setup and make it noticeably sharper.

\subsection{Keep \texttt{CLAUDE.md} alive}

\begin{tipbox}
Treat \verb|CLAUDE.md| as the single source of truth for conventions Claude should follow. Every time you correct Claude twice on the same thing (``use \verb|\citet| not \verb|\cite|'', ``British English'', ``do not rewrite section~2''), add a one-line rule. You are training future sessions.
\end{tipbox}

\subsection{Modularise with \texttt{\textbackslash input}}

One file per section makes diffs small, reviews fast, and lets Claude edit one section without touching others. A good default:

\begin{lstlisting}[style=tex]
main.tex
sections/
  introduction.tex
  background.tex
  methods.tex
  results.tex
  discussion.tex
  conclusion.tex
references.bib
figures/
\end{lstlisting}

\subsection{Use \texttt{/review} before pushing}

\verb|/review| is Claude Code's built-in slash command that runs a critical second pass over your uncommitted changes. Running it before \verb|git push| catches over-eager edits, invented citations, and broken cross-references before they reach Overleaf.

\subsection{Build verification in CI}

Add a GitHub Action that runs \verb|latexmk -pdf main.tex| on every push. Two benefits:

\begin{enumerate}[leftmargin=*]
  \item You find out about compile failures immediately, not when your co-author pulls.
  \item Claude can read the CI log after a failed build and self-correct next time.
\end{enumerate}

A minimal workflow:

\begin{lstlisting}[style=shell]
# .github/workflows/build.yml
name: Build PDF
on: [push, pull_request]
jobs:
  build:
    runs-on: ubuntu-latest
    steps:
      - uses: actions/checkout@v4
      - uses: xu-cheng/latex-action@v3
        with:
          root_file: main.tex
      - uses: actions/upload-artifact@v4
        with:
          name: paper
          path: main.pdf
\end{lstlisting}

\subsection{\texttt{latexdiff} for co-author-friendly revisions}

When you ship a revision, co-authors want to see what changed without reading the whole paper. \verb|latexdiff| produces a change-bar PDF between two commits:

\begin{lstlisting}[style=shell]
$ git show v1:main.tex > old.tex
$ latexdiff old.tex main.tex > diff.tex
$ latexmk -pdf diff.tex
\end{lstlisting}

\noindent Ask Claude to automate this as a script in the repo.

\subsection{One branch per revision round}

Scoped branch names --- for example \verb|claude/shorten-intro|, \verb|claude/reviewer-response|, or \verb|claude/notation-cleanup| --- keep work reversible and easy to review. Merge via pull request when you are happy; you get free code review on prose.

\subsection{Let Claude write your commit messages}

Over a paper's lifetime this becomes a per-change log: ``Tighten intro to fit R1 word limit'', ``Standardise notation for learning rate'', ``Add table~3 (ablation)''. Invaluable when a co-author asks ``what did we change last week?''

\subsection{Hooks: lint on save}

If you have \verb|chktex| installed, a \verb|PostToolUse| hook in \verb|.claude/settings.json| can run it automatically on every edited \verb|.tex| file:

\begin{lstlisting}[style=shell]
{
  "hooks": {
    "PostToolUse": [
      {
        "matcher": "Edit|Write",
        "hooks": [
          { "type": "command", "command": "chktex -q \"$CLAUDE_FILE_PATH\" || true" }
        ]
      }
    ]
  }
}
\end{lstlisting}

\noindent Claude sees the lint output and can correct on the spot.

\subsection{Ignore build artefacts aggressively}

Every file Overleaf generates locally (\verb|.aux|, \verb|.log|, \verb|.bbl|, \verb|.fls|, \verb|.fdb_latexmk|, \verb|.synctex.gz|) should be in \verb|.gitignore|. Syncing these back through Overleaf causes noisy commits and occasional merge grief.

\subsection{Use plan mode for big edits}

For anything larger than a section, ask Claude to plan before it writes:

\begin{lstlisting}[style=prompt]
> Plan how you would restructure the results section to lead with the headline finding. Do not edit yet; write the plan to plan.md.
\end{lstlisting}

\noindent Read the plan, push back on what you disagree with, then approve.

\subsection{Ask for a report, not a rewrite}

Before a destructive edit, ask Claude to produce a list. ``Which \verb|\cite| keys are unused?'', ``Which sentences are longer than 40 words?'', ``Which figures have no caption?''. You then decide what to act on.

\subsection{Use custom slash commands}

Repeated tasks become slash commands in \verb|.claude/commands/|. A \verb|/tighten| command that shortens the current section, or a \verb|/ack| that drafts an acknowledgements paragraph, saves retyping the prompt every time.
